\begin{qfdefinition}[Fixed-income market structure: segments, trading protocols, venues and dealer risk-bearing, kept apart (U.S., as surveyed in 2019)]\label{def:fixed-income-market-structure}
\qffield{Statement} Except for a few segments, fixed-income instruments trade in dealer-oriented over-the-counter (OTC) search markets that differ significantly from equity markets \citep[sec.~I, p.~1]{wide-bsv-fi-survey}. Three things that the word ``market'' runs together are kept apart here.
\begin{itemize}
\item \emph{Segments.} U.S. Treasury securities are ``in many ways the simplest'' fixed-income instruments: low default risk, high volume, standardised contracts and generally liquid markets (p.~1). Corporate bond trading is fragmented because each issue of an issuer is a distinct contract, not fungible with the others; individual issues did not trade on 52 percent of days in Edwards, Harris and Piwowar's sample, and in 2016 9.3 percent of TRACE bonds never traded and 1.5 percent traded every day (pp.~5 and~14). Structured products split between to-be-announced (TBA) agency MBS, forward contracts on any deliverable pool, the most actively traded structured securities, and specified pools and other structured bonds, where most MBS and ABS did not trade at all in one 21-month sample (pp.~5 and~13--14). Municipal bond trading is also highly fragmented, with over 1.5 million distinct bonds (p.~16).
\item \emph{Trading protocols}, the rules by which a client obtains a price. Dealers broadcast indicative bids and offers (``runs''); clients negotiate by phone or chat, or send bid- and offer-wanted lists; in an electronic request for quote (RFQ), as Hendershott and Madhavan describe it on MarketAxess, the client chooses how many dealers to query for a given size and a deadline, sees all quotes at the deadline and may trade with the best (pp.~2--3 and~7--8). Fixed-income quotes are typically indicative, not firm (p.~3).
\item \emph{Venues}, the systems on which trading happens: interdealer brokers, about 70 percent of U.S. Treasury interdealer trading being on electronic limit-order-book platforms where principal trading firms also take part; RFQ platforms for customers; and newer all-to-all venues on which institutions trade with each other (pp.~8 and~11).
\item \emph{Dealer risk-bearing.} Outside Treasuries and TBA MBS, trading is dominated by customer--dealer OTC trades, traditionally on a principal basis, the dealer absorbing the order into inventory. Dealers have reduced capital commitment in corporate bonds, relying more on prearranged trades in which the dealer has already located the other side, and on shorter intermediation chains; nonbank dealers' share of customer--dealer corporate bond volume rose from 2.4 percent before the 2007--2009 crisis to 12.5 percent in 2014--2016 (p.~4).
\end{itemize}

\qffield{Conditions} The survey describes U.S. markets with data to about 2017--2018; its electronic-share figures come from industry reports: an industry report published in July 2016 estimated that 19 percent of U.S. investment-grade corporate bond trading was electronically facilitated, mostly by RFQ; the observation period behind that estimate is not given in the passage (p.~2). Hendershott and Madhavan find that periodic one-sided electronic auctions are ``a viable and important source of liquidity even in inactively traded instruments'', a compromise between bilateral search and continuous double auctions \citep[abstract]{wide-hendershott-madhavan}; in the survey's account RFQ use is greatest for recently issued, investment-grade, large-issue bonds (p.~3). O'Hara and Zhou, combining MarketAxess and TRACE data, find electronic RFQ trading ``fairly small and segmented'' but with wide-ranging effects on transaction costs and execution quality in electronic and voice trading and the interdealer market \citep[abstract]{wide-ohara-zhou-electronic}. Both are abstract-level checks.

\qffield{Practice} Interpretation: a dealer's balance sheet is not a venue; the dealer's willingness to hold inventory is a source of liquidity, reached through a protocol (voice, RFQ) on a venue or bilaterally. The survey links the reduction in dealers' capital commitment to post-crisis capital requirements and the Volcker Rule, and notes that counting only bank-affiliated dealers' capital may overstate the effect on markets because nonbank dealers, electronic intermediation and customers' liquidity provision through prearranged trades respond (pp.~4 and~32--33). The generic inventory mechanics are in \hyperref[def:dealer-inventory-risk]{dealer inventory risk} and \hyperref[meth:dealer-quote-adjustment]{dealer quote adjustment}, from equity and floor examples.

\qffield{Pitfalls} That government bonds are the most standardised and electronically tradable holds for U.S. Treasuries in the survey, whose electronic share is interdealer; customer-to-dealer Treasury trading relies on electronic communication for a large share but was not measured as electronic execution (p.~3). That standard credit indices trade more centrally than bespoke risk, and that swaps often remain dealer-mediated because contracts differ, are not measured here: the survey mentions only that most standardised derivatives, including interest rate swaps, must now be centrally cleared, and a study linking post-trade transparency to lower CDS trading costs (pp.~30 and~32).

\emph{U.S. rules that separate the cases}.
\begin{itemize}
\item \emph{Clearing.} The CFTC's clearing requirement covers named classes of fixed-to-floating, basis, FRA and overnight-index swaps by currency, index and maturity range (for example SOFR and Fed Funds OIS), and the North American and European \emph{untranched CDS indices}: CDX.NA.IG and HY, and iTraxx Europe, Crossover and HiVol, for stated tenors and series. No single-name CDS class is listed \citep[17 CFR 50.4(a)--(b), last amended 91~FR 57086, 8 September 2026]{ecfr-17cfr-cftc-swaps}.
\item \emph{Execution.} A swap subject to the trade-execution requirement, once a facility has made it ``available to trade'' by a self-certified determination weighing buyers and sellers, frequency, volume, participants, spreads and resting quotes (37.10), must be executed on a swap execution facility. Unless it is a block trade, it must go through an order book or through a request-for-quote system that sends the request ``to no less than three market participants'' who are not affiliates, operating with an order book (37.9(a)). Other swaps may trade by any method the facility offers (37.9(c)).
\end{itemize}
So the rules, not only the contracts, separate standard indices from single names and bespoke risk. \emph{Which swaps are available to trade}. The CFTC's table of MAT filings lists the determinations of October 2013 to March 2014 (Javelin, trueEX, TW SEF, MarketAxess, Bloomberg), for certain interest-rate swaps and CDS. It also lists TW SEF's voluntary submission of 22 May 2023, approved on 7 July 2023 \citep[``Swaps Made Available to Trade Determination'' table]{cftc-smatd-filings}. That submission recalls its 2013 determination: USD and GBP LIBOR fixed-to-floating swaps in the 2- to 30-year benchmark tenors. It sets out that the CFTC's August 2022 amendment to 50.4 moved the clearing requirement from LIBOR swaps to SONIA and SOFR overnight index swaps of 7 days to 50 years, and it proposes MAT status for SOFR and SONIA OIS \citep[TW SEF letter of 22 May 2023, pp.~1--3]{twsef-mat-2023}.

It also notes that ``on an RFQ platform like TW SEF, there are no resting bids and offers'', and measures liquidity there by the number of providers streaming quotes. That is one operator's description of its own protocol. The share of swap trading done by dealer-to-client RFQ is not measured in these sources. RFQ does not dominate U.S. investment-grade corporate bonds in the survey's figure: 19 percent was electronically facilitated, mostly by RFQ, in the estimate of a July 2016 industry report (p.~2); RFQ is one protocol among voice negotiation, runs and lists.

That products trade by RFQ because they are less standardised and active is not a finding: the survey treats why fixed-income trading remains dealer-oriented and OTC as an open question (pp.~37--38). One corporate-bond sample does not describe every fixed-income market.

\qffield{See also} \hyperref[rule:us-eu-bond-transparency]{bond transparency rules}; \hyperref[ev:bond-transparency-trading-costs]{bond transparency and trading costs}; \hyperref[def:dealer-inventory-risk]{dealer inventory risk}; \hyperref[mech:etf-creation-redemption]{ETF creation and redemption} (corporate bond ETFs and dealer inventories); \hyperref[def:portfolio-trading]{portfolio trading}.

\end{qfdefinition}

\begin{qfrule}[Bond transparency rules: U.S. TRACE reporting and dissemination, and EU post-trade deferrals (as of October 2026)]\label{rule:us-eu-bond-transparency}
\qffield{Statement} Transparency rules separate \emph{reporting} to a regulator or its facility from \emph{public dissemination}, and \emph{post-trade} publication of completed trades from \emph{pre-trade} publication of quotes or orders.
\begin{itemize}
\item \emph{U.S. reporting.} Under FINRA Rule 6730 as displayed on 15 October 2026, each member that is a party to a transaction in a TRACE-eligible security must report it ``as soon as practicable, but no later than within 15 minutes of the Time of Execution'', with stated exceptions, some of which allow 60 minutes \citep[Rule 6730(a)]{finra-rule-6730}.
\item \emph{U.S. dissemination.} Under Rule 6750, FINRA disseminates information on all transactions in TRACE-eligible securities, Rule 144A transactions included, immediately upon receipt, except that certain collateralised mortgage obligation trades are disseminated in aggregate weekly and monthly; individual trades in on-the-run nominal coupon Treasuries are disseminated at the end of the day; and other Treasuries, foreign sovereign debt, CMBS, CDOs, some CMO trades and certain affiliate, list-offering and takedown transactions are not disseminated, FINRA being allowed to publish aggregated statistics \citep[Rule 6750(a)--(d) and .01]{finra-rule-6750}.
\item \emph{EU post-trade deferrals.} ESMA's manual of 12 May 2026 gives worked examples of the bond regime that applies from 2 March 2026 under amended MiFIR Article 11 and RTS~2 Article 8a. For a trade in a group-1 (most liquid sovereign) bond on a trading venue, a medium-size trade in category~1 is published within 15 minutes, and a large one in category~3 has its price published within one trading day and its volume within one week; competent authorities may add supplementary deferrals for sovereign bonds, omitting volume for up to six months beyond the initial deferral, or publishing weekly aggregates \citep[secs.~4.2.8 and~4.2.8.2, pp.~262--272]{esma-transparency-manual-2026}.
\end{itemize}

\qffield{Conditions} Dates and status:
\begin{itemize}
\item \emph{History, United States.} TRACE reporting of corporate bonds began in July 2002 within 75 minutes, with immediate public dissemination only for about 500 large investment-grade and 50 high-yield bonds; dissemination was widened in 2003--2005, and the reporting lag cut to 45 (October 2003), 30 (October 2004) and 15 minutes (July 2005), with all registered corporate bonds disseminated by January 2006. ABS and MBS reporting began in May 2011 without dissemination, which was phased in for TBAs (2012), specified pools (2013), Rule 144A bonds (2014), ABS (2015) and CMOs (2017). Treasury reporting began in July 2017, with no public dissemination at the time \citep[Table~2, p.~9]{wide-bsv-fi-survey}. Disseminated sizes were capped, for example ``5MM+'' for investment-grade and ``1MM+'' for high-yield corporate trades (p.~10, footnote~10); current caps may differ.
\item \emph{Changes and proposals, United States.} FINRA's filing SR-FINRA-2024-004 proposed reducing the 15-minute reporting limit to one minute, with exceptions for firms with de minimis activity and for manual trades; SR-FINRA-2025-008, by its title, maintains ``the currently effective 15-minute outer limit'' for the securities covered by the earlier filing, and the rulebook lists both as effective 8 June 2026 \citep[filing pages]{finra-sr-2024-004}. SR-FINRA-2026-009, approved by the SEC on 24 July 2026, expands the indicator that suppresses from dissemination qualifying same-day, same-price, same-security transactions between a member and its \emph{non-member} affiliate to transactions between \emph{member} affiliates, renaming it the Affiliate--Principal Transaction indicator; FINRA said it would set an implementation date at least 12 months after approval and allow voluntary earlier use \citep[pp.~47909--47911]{sec-34-105991}. The rulebook lists it as effective 2 August 2027.

Until then the presently effective exclusion is the non-member-affiliate one; the displayed text of Rule 6750(d)(1), which already reads ``Affiliate---principal transaction indicator'', is the future wording. The proposal of the SEC's Fixed Income Market Structure Advisory Committee to delay dissemination of very large trades by 48 hours was a proposal, not a rule (survey, p.~29).
\item \emph{European Union.} Fixed-income securities traded without pre- or post-trade transparency in Europe until MiFID~II applied in January 2018, which required reporting of bond transactions to a national regulator and real-time publication for liquid, non-block trades, and pre-trade publication of bids and offers for small trades in a limited set of about 200 sovereign bonds (survey, p.~30; historical). On 10 October 2025 ESMA stated that the Commission's amending RTS~1 and RTS~2 would apply most revised bond transparency requirements from 2 March 2026 following publication in the Official Journal, which had not yet happened, and invited market participants to anticipate them \citep[p.~4]{esma-mifir-review-statement-2025}. \emph{Level~1 and RTS~2 as amended}.
\begin{itemize}
\item \emph{Pre-trade.} Regulation (EU) 2024/791 confines venue pre-trade transparency for bonds, structured finance products and emission allowances to central limit order books and periodic auctions. Venues using those systems ``shall make public current bid and offer prices and the depth of trading interests'' continuously during trading hours, calibrated by trading system \citep[Art.~1, point~(5), replacing MiFIR Art.~8(1)--(2) and deleting Art.~8(3)--(4)]{eu-2024-791}. Request-for-quote and voice systems are thus no longer named for bonds.
\item \emph{Deferral categories.} Bond deferrals are organised in five categories: medium size liquid (1), medium size illiquid (2), large liquid (3), large illiquid (4) and very large (5). When a deferral lapses, ``all the details of the transactions on an individual basis shall be published''. Sovereign-debt supplementary deferrals are limited to the issuing Member State's instruments, with volume omitted for up to six months (Art.~1, point~(9), amending MiFIR Art.~11(1a) and~(3)).
\item \emph{Maximum durations.} The amending RTS~2 sets the maxima \citep[Art.~1, point~(7), inserting Art.~8a(1) into Delegated Regulation 2017/583]{eu-2025-1246}: 15 minutes for price and volume in category~1; end of the trading day in category~2; price by the end of the next trading day and volume within one week in category~3; price by the end of the second trading day and volume within two weeks in category~4; four weeks for both in category~5. It also requires post-trade publication ``within five minutes'' (amended Art.~7(4)). It was published in the Official Journal on 3 November 2025 (OJ~L 2025/1246), and its Article~1 applies from 2 March 2026. The size and liquidity thresholds of Annex~III, Table~2.6, were not transcribed.
\end{itemize}
These agree with ESMA's worked examples above.
\item \emph{United Kingdom}. The FCA's PS24/14 (November 2024) replaces the inherited regime, whose ``overly lengthy publication deferrals'' it criticises, with four size groups for bonds: real-time publication up to the first threshold, then deferrals of price and size together to the end of the trading day after execution, two weeks, and three months for the largest trades; it considered publishing price after four weeks and volume after three months for the largest group but chose ``both after three months'' (PDF pp.~42--43); most of its rules take effect on 1 December 2025 (PDF p.~14) \citep[PS24/14, chaps.~2--5, PDF pp.~12, 14 and~42--43]{fca-ps24-14}. The thresholds by issuer, issue size and maturity (Table~8) and the Handbook text (MAR~11) were not transcribed.
\end{itemize}

\qffield{Practice} Interpretation: a TRACE report is a regulatory obligation of the dealer; what the public sees is set by Rule 6750, so the reported and the disseminated sets of U.S. trades differ, most visibly for Treasuries. In the EU the deferral category, not only the instrument, decides when price and volume become public. Trade-size caps and deferrals let dealers ``hedge or sell a portion of [their] position before disclosure'', the reason the survey gives for U.S. regulators' initial delays (p.~29).

\qffield{Pitfalls} TRACE is FINRA's reporting facility for all TRACE-eligible securities, of which corporate bonds are one class; reports go to FINRA, and public dissemination is a separate step with exceptions. Whether TRACE improved post-trade transparency is an evidence question (\hyperref[ev:bond-transparency-trading-costs]{bond transparency and trading costs}). No EU statement of aims supports the view that European transparency rules aim to make fixed-income liquidity easier to measure without forcing all products into a CLOB. The 2002--2017 history is from a survey, not from the rule filings themselves.

\qffield{See also} \hyperref[def:fixed-income-market-structure]{fixed-income market structure}; \hyperref[ev:bond-transparency-trading-costs]{bond transparency and trading costs}; \hyperref[def:best-execution]{best execution}.

\end{qfrule}

\begin{qfevidence}[Bond transparency and trading costs: U.S. post-trade transparency evidence, the open question of large trades, and what data allow clients to measure]\label{ev:bond-transparency-trading-costs}
\qffield{Statement} Edwards, Harris and Piwowar estimate average transaction costs by trade size for every U.S. corporate bond that traded more than nine times between January 2003 and January 2005, from a complete record of OTC secondary trades. Costs fall significantly with trade size; highly rated, recently issued and near-maturity bonds cost less; costs are lower for bonds with transparent trade prices and drop when TRACE starts to disseminate their prices; and the results ``suggest that public traders benefit significantly from price transparency'' \citep[abstract]{wide-edwards-harris-piwowar}. In the survey's account, their estimates fall from 75 basis points for \$5{,}000 trades to under 10 basis points for \$1~million or more, the reverse of the equity pattern \citep[p.~5]{wide-bsv-fi-survey}.

Three studies with different designs, a stratified random sample of BBB bonds phased in in April 2003 (Goldstein, Hotchkiss and Sirri), the disseminated-against-undisseminated comparison of Edwards, Harris and Piwowar, and a difference-in-differences on insurance-company trades (Bessembinder, Maxwell and Venkataraman), all conclude that public reporting through TRACE is associated with significant reductions in customer execution costs; the last also suggests a liquidity externality to bonds not yet reported. Similar associations are reported for CDS, TBA MBS and Rule 144A bonds (survey, p.~30).

\qffield{Conditions} U.S. evidence from 2002--2017 samples; Rule 2232 as displayed on 15 October 2026; other regulatory facts as of the survey (2019). These are associations in U.S. corporate bond, CDS, TBA and Rule 144A samples from the phase-in periods, as summarised by a survey. The survey states the limits: much of the evidence concerns execution costs of completed trades, which may reflect order splitting and do not capture the cost and difficulty of executing large orders; the optimal degree of post-trade transparency ``has not been established''; dealers, and in some cases large buy-side investors, are often critical of greater transparency; and the regime optimal for one fixed-income market may not be optimal for another (pp.~38--39).

Back, Liu and Teguia's model, as the survey reports it, has trade disclosure conveying the dealer's information and reducing its rents when it disposes of inventory (p.~29), a model prediction. Pre-trade transparency in bonds is a matter of open questions and proposals in the survey, such as the public display of customer interest advocated by Harris (2015), not of evidence (p.~37).

\qffield{Practice} What data let clients measure, as of the survey: TRACE and MSRB transaction data; an enhanced TRACE file for researchers with masked dealer identities and actual sizes; aggregators that combine trades, dealer runs and platform quotes for buy-side institutions (pp.~7--8 and~10). Fixed-income markets have no analogue of SEC Rules 605 and 606 on execution quality and order routing, and without consolidated quotations many participants cannot see the best available prices or identify trade-throughs; FINRA and the MSRB adopted best-execution rules (pp.~31--32).

Since 14 May 2018 FINRA Rule 2232(c) requires a confirmation to show the mark-up or mark-down when a member trades as principal in a corporate or agency debt security with a non-institutional customer and makes offsetting trades of at least that size on the same trading day, looking through non-arm's-length affiliate trades; (d) exempts functionally separate trading desks and same-day fixed-price-offering sales \citep[paras.~(c)--(f)]{finra-rule-2232}; \citep[effective date and MSRB paragraph]{finra-rn-17-08}. FINRA's notice reports a parallel MSRB rule for municipal securities, with the same effective date. Measure costs by trade size, rating and segment, and state which trades are visible: completed trades only, capped sizes, deferred or undisseminated classes.

\qffield{Pitfalls} Publishing pre-trade interest in illiquid bonds can reveal inventory needs and increase dealer risk, so policy balances monitoring against exposing liquidity. The sources give this as dealers' argument and an open question, and as the survey's reason for U.S. reporting delays; none measures the effect of pre-trade transparency on bond liquidity. That electronification changes execution quality and the governance of trading desks is not established causally: O'Hara and Zhou report effects of electronic RFQ trading in their abstract, and nothing here concerns desk governance. Transparency is not shown here to help or to harm liquidity in general.

\qffield{See also} \hyperref[rule:us-eu-bond-transparency]{bond transparency rules}; \hyperref[def:fixed-income-market-structure]{fixed-income market structure}; \hyperref[def:spread-measures]{spread measures}; \hyperref[def:best-execution]{best execution}; \hyperref[meth:transaction-cost-analysis]{transaction cost analysis}.

\end{qfevidence}
